% !TeX program = xelatex
%
% UFAL Presentation Template — college-beamer theme with UFAL branding.
% Replace the title/author/date below and the instructional text in each
% frame with your own content. See README.md for details.
%
\documentclass[9pt]{beamer}
\usepackage[szu,en]{collegeBeamer}

% `fontspec` (used below for the \cascadia monospace font in the author's
% email) only works under XeTeX/LuaTeX -- it throws a fatal error under
% pdfTeX. Some setups (e.g. Overleaf, if the compiler is set to pdfLaTeX
% and ignores the "% !TeX program = xelatex" magic comment on line 1) still
% compile this file with pdflatex, so \cascadia is defined conditionally:
% a real custom font under XeTeX/LuaTeX, a plain \ttfamily fallback under
% pdfTeX (which also needs inputenc/fontenc for the accented characters
% used elsewhere, e.g. "Maceió", since pdfTeX is not UTF-8-native like
% XeTeX/LuaTeX).
\usepackage{iftex}
\ifxetex
    \usepackage{fontspec}
    \IfFontExistsTF{Cascadia Code}
        {\newfontfamily\cascadia{Cascadia Code}}
        {\newfontfamily\cascadia{DejaVu Sans Mono}}
\else
    \ifluatex
        \usepackage{fontspec}
        \IfFontExistsTF{Cascadia Code}
            {\newfontfamily\cascadia{Cascadia Code}}
            {\newfontfamily\cascadia{DejaVu Sans Mono}}
    \else
        \usepackage[utf8]{inputenc}
        \usepackage[T1]{fontenc}
        \newcommand{\cascadia}{\ttfamily}
    \fi
\fi
\usetikzlibrary{positioning,arrows.meta}

% biblatex: edit referencias.bib with your own references
\usepackage[backend=biber,style=numeric,sorting=none]{biblatex}
\addbibresource{referencias.bib}

\renewcommand{\qaStr}{\textsl{Thank you for your attention! \\ Questions and comments are welcome.}}

\definecolor{hrefcol}{RGB}{55, 50, 204} % Example: blue color

% meta-data -- replace with your own
\title{
{\large\color{blue!60!black}\textbf{Your Department or Field}}\\[0.8em]
Your Presentation Title Goes Here
}
% the theme's title-page template always appends "\\" after \insertsubtitle;
% an empty subtitle makes that "\\" have nothing to end, causing a LaTeX error,
% so a near-invisible placeholder (thin space) is used instead of leaving it undefined
\subtitle{\,}

\author{
Your Name\\
{\footnotesize\cascadia\href{mailto:your.email@example.com}{your.email@example.com}}}

\date{
Maceió, Brazil \\
Month Year
}

% document body
\begin{document}

    \maketitle

    \section{Introduction}

    \begin{frame}{Context and Motivation}
        \begin{itemize}
            \item Describe the problem domain and why it matters \cite{example2024}
            \item State the current situation or practice
            \item Point out the limitations of existing approaches
            \item Introduce, in one line, the direction your work takes
        \end{itemize}
    \end{frame}

    \begin{frame}{Background}
        \small
        \begin{itemize}
            \itemsep1pt
            \item Key concept 1 your audience needs before the next slides \cite{example2024}
            \item Key concept 2
            \item Key concept 3
            \item How these concepts relate to each other
        \end{itemize}
        \vskip0.3em
        \centering\footnotesize\ctextbf{One-sentence statement of the gap your work addresses}
    \end{frame}

    \begin{frame}{General Objective}
        State, in one or two sentences, the overall goal of the work.

        \vskip1em
        \begin{block}{Research Questions}
            \begin{enumerate}
                \item First research question
                \item Second research question
                \item Third research question
            \end{enumerate}
        \end{block}
    \end{frame}

    \begin{frame}{Specific Objectives}
        \begin{itemize}
            \item First specific objective
            \item Second specific objective
            \item Third specific objective
            \item Fourth specific objective
            \item Fifth specific objective
        \end{itemize}
    \end{frame}

    \section{Related Work}

    \begin{frame}{Why This Work Matters}
        \small
        \begin{itemize}
            \itemsep1pt
            \item What the closest related work already does \cite{example2024}
            \item What it does not cover
            \item How your work is positioned relative to it
        \end{itemize}
        \vskip0.2em
        \begin{block}{Gaps Addressed}
            \begin{enumerate}
                \itemsep0pt
                \item First gap your work fills
                \item Second gap your work fills
            \end{enumerate}
        \end{block}
    \end{frame}

    \section{Methodology}

    \begin{frame}{Pipeline Overview}
        \begin{center}
        \resizebox{\textwidth}{!}{%
        \begin{tikzpicture}[node distance=0.5cm, every node/.style={font=\scriptsize}]
            \tikzset{box/.style={draw=maincolor, thick, rounded corners, fill=morelightgray, text width=2.4cm, align=center, minimum height=1.5cm}}
            \node[box] (a) {Input / data};
            \node[box, right=of a] (b) {Processing step 1};
            \node[box, right=of b] (c) {Processing step 2};
            \node[box, right=of c] (d) {Processing step 3};
            \node[box, right=of d] (e) {Output / evaluation};
            \foreach \x/\y in {a/b, b/c, c/d, d/e}{\draw[-{Latex}, thick] (\x) -- (\y);}
        \end{tikzpicture}%
        }
        \end{center}
        \vskip0.5em
        Explain, in one or two sentences, the pipeline shown above.
    \end{frame}

    \begin{frame}{Data and Setup}
        \begin{itemize}
            \item Where the data comes from and how it was collected
            \item Size and relevant characteristics of the dataset
            \item Key configuration choices, parameters, or constraints
            \item Ethical or privacy considerations, if applicable
        \end{itemize}
    \end{frame}

    \begin{frame}{Experimental Design}
        \begin{center}
        \small
        \renewcommand{\arraystretch}{1.15}
        \begin{tabular}{|c|c|c|}
            \hline
            \textbf{Condition} & \textbf{Variable A} & \textbf{Variable B} \\
            \hline
            Group 1 & value & value \\
            \hline
            Group 2 & value & value \\
            \hline
            Group 3 & value & value \\
            \hline
        \end{tabular}
        \end{center}
        \vskip0.5em
        Briefly explain the comparison design and the evaluation metrics used.
    \end{frame}

    \section{Results}

    \begin{frame}{Main Finding}
        \begin{itemize}
            \item State the headline result plainly
            \item Give the supporting number or comparison
            \item Note what this result does -- and does not -- show
        \end{itemize}
        \vskip0.3em
        \centering\footnotesize\ctextbf{One-sentence takeaway for this slide}
    \end{frame}

    \begin{frame}{Detailed Results}
        \begin{itemize}
            \item Breakdown point 1, e.g.\ per group/class/condition
            \item Breakdown point 2
            \item Breakdown point 3
            \item Statistical test used and its outcome \cite{example2024}
        \end{itemize}
    \end{frame}

    \begin{frame}{Qualitative Observations}
        \small
        \begin{itemize}
            \item A concrete example or case illustrating a pattern
            \item Another concrete example
            \item What these examples suggest
        \end{itemize}
    \end{frame}

    \section{Concluding Remarks}

    \begin{frame}{Answer to the Research Question}
        \begin{itemize}
            \item How your results answer the first research question
            \item How your results answer the second research question
        \end{itemize}
        \vskip0.3em
        \centering\footnotesize\ctextbf{Practical implication of your findings}
    \end{frame}

    \begin{frame}{Contributions}
        \small
        \begin{itemize}
            \item First contribution
            \item Second contribution
            \item Third contribution
            \item Fourth contribution
        \end{itemize}
    \end{frame}

    \begin{frame}{Limitations \& Future Work}
        \small
        \begin{itemize}
            \item First limitation
            \item Second limitation
        \end{itemize}
        \vskip0.3em
        \textbf{Future work:} list two or three concrete next steps here.
    \end{frame}

    \bibliographpage

    \QApage

\end{document}
